% !TEX program = xelatex
\documentclass[a4paper]{article}
\usepackage{amsthm}
\usepackage{amssymb}
\usepackage{bm}
\usepackage{mathtools}
\usepackage[x11names]{xcolor}
\usepackage{xparse}
\usepackage{fontspec}
\usepackage{unicode-math}
\setromanfont{DovesType-Regular.otf}
\setsansfont{Andika}
\setmathfont{Asana Math}[Scale=1]

\usepackage{varwidth}
\usepackage{siunitx}
\usepackage{graphicx}
\usepackage[margin=1cm]{geometry}
\usepackage[most]{tcolorbox}
\usepackage{pgfplots}
\pgfplotsset{compat=newest}
\tcbuselibrary{skins,xparse,poster,breakable}
% \usetikzlibrary{fadings}
\usetikzlibrary{calc, plotmarks, shapes, shapes.geometric, positioning, angles, intersections, quotes, through, patterns, turtle, arrows.meta}
\usetikzlibrary{decorations.markings}
% \usepackage{etoolbox}
\usepackage{tkz-euclide}

%%%%%%%%%%%%%%%%%%%%%%%%%%%%%%%%%%%%%%%%%%%%%%%%%%%%%%%%%
\newcommand\markangle[9]{% origin X Y radius radiusmark mark colour opacity
%  % fill red circle offset-from-centre
  \begin{scope}
    \path[clip] (#1) -- (#2) -- (#3);
    \fill[color=#7,fill opacity=#8,draw=black,name path global=pcircle]  % global declaration required otherwise pcircle is not seen by the `named intersections=' lines below.
    (#1) circle (#4);
  \end{scope}
  % middle calculation
  \path[name path=line one] (#1) -- (#2);
  \path[name path=line two] (#1) -- (#3);
  \path[%
  name intersections={of=line one and pcircle, by={inter one}},
  name intersections={of=line two and pcircle, by={inter two}}
  ] (inter one) -- (inter two) coordinate[pos=#9] (place);
  % put mark
  \node at ($(#1)!#5!(place)$) {\scriptsize{#6}};
}

% \newcommand{\condSoln}[2]{\ifcsdef{r@#1}{#2}{}}

% \newcommand\fadingtext[3][]{%
%    \begin{tikzfadingfrompicture}[name=fading letter]
%      \node[text=transparent!0,inner xsep=0pt,outer xsep=0pt,#1] {#3};
%    \end{tikzfadingfrompicture}%
%    \begin{tikzpicture}[baseline=(textnode.base)]
%      \node[inner sep=0pt,outer sep=0pt,#1](textnode){\phantom{#3}};
%      \shade[path fading=fading letter,#2,fit fading=false]
%      (textnode.south west) rectangle (textnode.north east);%
%    \end{tikzpicture}%
% }

\definecolor{JISpurple}{RGB}{89,72,122}
\definecolor{JISivory}{RGB}{241,234,221}
\definecolor{JIStaupe}{RGB}{183,156,154}
\definecolor{PaleGreen}{RGB}{240,255,240} % 'Honeydew'

\AddToHook{shipout/background}{%
    \put (0in,-\paperheight){\includegraphics[width=\paperwidth,height=\paperheight]{images/R10V5-R.png}}%
}

\newcommand\numberthis{\addtocounter{equation}{1}\tag{\theequation}}

\newtcolorbox{MyOuterBox}{%
  enhanced,
  breakable,
  frame style=JISpurple,
  colback=JISivory,
  colframe=JISpurple,
  title={\includegraphics[width=0.9cm,height=0.9cm]{images/JIS Final Logo FA-02.png}\raisebox{3mm}{\Large{Maths Challenge}\hspace{26em} \Large{\bfseries\sffamily 5}}},
}

\newtcolorbox{MyInnerBox}[2][]{enhanced,%empty,
coltitle=JISpurple,colback=white,
fonttitle=\bfseries\sffamily,
attach boxed title to top left={yshift=-1.5mm},
boxed title style={empty, size=small, top=1mm, bottom=0pt},
varwidth boxed title=0.5\linewidth,
frame code={
  \path (title.east|-frame.north) coordinate (aux);
\path[draw=JISpurple, line width=0.5mm, rounded corners,fill=white]
(frame.west) |- ([xshift=-2.5mm]title.north east) to[out=0, in=180] ([xshift=7.5mm]aux)-|(frame.east)|-(frame.south)-|cycle;
},
title={#2},#1}

\newtcolorbox{MyInnerSplitBox}[2][]{enhanced,%empty,
bicolor,sidebyside,sidebyside align=top seam,
righthand width=5.5cm,colbacklower=white,
coltitle=JISpurple,colback=white,
fonttitle=\bfseries\sffamily,
attach boxed title to top left={yshift=-1.5mm},
boxed title style={empty, size=small, top=1mm, bottom=0pt},
varwidth boxed title=0.5\linewidth,
frame code={
  \path (title.east|-frame.north) coordinate (aux);
\path[draw=JISpurple, line width=0.5mm, rounded corners,fill=white]
(frame.west) |- ([xshift=-2.5mm]title.north east) to[out=0, in=180] ([xshift=7.5mm]aux)-|(frame.east)|-(frame.south)-|cycle;
},
title={#2},#1}


\newtcolorbox{MySolutionBox}{%
  enhanced,
  breakable,
  frame style=JISpurple,
  colback=PaleGreen, colframe=green,
  title={\Large Solution},
  drop fuzzy shadow,
  halign=left,
}

%%%%%%%%%%%%%%%%%%%%%%%%%%%%%%%%%%%%%%%%%%%%%%%%%%
\newtoggle{SOLUTION}
%%% Uncomment the appropriate line below to show solutions %%%
% \toggletrue{SOLUTION}
\togglefalse{SOLUTION}
%%%%%%%%%%%%%%%%%%%%%%%%%%%%%%%%%%%%%%%%%%%%%%%%%


%%%%%%%%%%%%%%%%%%%%%%%%%%%%%%%%%%%%%%%%%%%%%%%%%%
%%%%%%            DOCUMENT BEGINS           %%%%%%
%%%%%%%%%%%%%%%%%%%%%%%%%%%%%%%%%%%%%%%%%%%%%%%%%%
\begin{document}


  \begin{MyOuterBox}
    \iftoggle{SOLUTION}{Here are the full, or partial solutions.
    }{
      Welcome to this week's Maths Challenge!\\
      Everyone is encouraged to take part.\\
      There are 3 to 5 merit points for correct solutions, depending on the difficulty of the problem and how impressed the marker is with your solution.\\
      There are house points for students who submit 7 or more correct solutions in one term.\\
      Solutions must be explained in detail, responses with just the answer may be ignored.\\
      Drop your solution in the box in the staffroom by Tuesday.
    }
       \begin{MyInnerBox}{Year 9 and below}
      This is a version of a very well known logic puzzle. Try to figure it out by yourself though!\par
      You are alone and lost, following a small path through a dark and eldritch forest. Shadows loom, strange rustles menace and the day is failing. You come to a fork in the path: there are two ancient oak trees and a peculiar sign grows in front of you and tells you that travelling down one path will lead you out of the forest, but walking down the other dooms you to a dreadful end. You may ask one question of one of the two wyrd oaks, they can only answer by pointing a gnarled branch. You are also told that one of the trees always lies while the other always tells the truth, but you have no way of telling which is which. What question should you ask to escape the deep dire wood?
      \iftoggle{SOLUTION}{%conditional output begin
      \begin{MySolutionBox}
        First, you have to realise that just asking a question like "Which is the way out of the forest?" will not work, for if you have asked the truth-teller, they will point to the correct path, if you have asked the liar it will point to the wrong path, but you have learned nothing because you do not know which oak is which.\par
        The trick is to smuggle two questions into one. If you ask one of the trees "Which path will the other tree point to if I ask it the correct way out?" let's see what happens.\par
        Case 1. You ask the lying oak.\par
        The lying oak knows the truthful tree will point to the way out if asked, but then lies and points to the path of doom.\par
        Case 2. You ask the truthful oak.\par
        The truthful tree knows the lying tree would point to the wrong path if asked for the way out, and tells the truth about that, thus pointing to the path of your gruesome demise.\par
        So when you ask this question, you simply take the path not pointed to, and you have won your freedom from the forest, you go on your way celebrating the defeat of the quirky quercus and the cleverness of you. Notice that you never find out which tree is the liar and which is the truth-teller.
      \end{MySolutionBox}
    }{}%conditional output end
    \end{MyInnerBox}


    \vspace{0.4cm}
          \begin{MyInnerBox}{Year 10 and above}
       If \[ \sin^{3}\theta + \cos^{3}\theta = \frac{11}{16}\] find the value of \[\sin\theta + \cos\theta\]\par
       Hint: Factorise the first equation. Let \(x = \sin\theta + \cos\theta\). You should obtain a cubic polynomial in \(x\). \(x = \frac{1}{2}\) is a root.
      \iftoggle{SOLUTION}{%conditional output begin
        \vspace{0cm}
      \begin{MySolutionBox}
        Useful identities to start with are \(a^{3}+b^{3} = (a+b)(a^{2}-ab+b^{3})\) and \(\sin^{2}\theta + \cos^{2}\theta = 1\). Using these we have,
        \begin{align*}
          \sin^{3}\theta + \cos^{3}\theta = \frac{11}{16} &= (\sin\theta + \cos\theta)(\sin^{2}\theta - \sin\theta\cos\theta + \cos^{2}\theta)\\
          \frac{11}{16}  &= (\sin\theta + \cos\theta)(1-\sin\theta\cos\theta) \numberthis \label{eq1}\\
          \text{Let } x &= \sin\theta + \cos\theta \implies x^{2} = \sin^{2}\theta + \cos^{2}\theta + 2\sin\theta\cos\theta\\
          x^{2} &= 1 + 2\sin\theta\cos\theta \implies \frac{x^{2}-1}{2} = \sin\theta\cos\theta\\
          \shortintertext{Substituting in \eqref{eq1}, we obtain a cubic in \(x\):}
          \frac{11}{16} &= x\left(1-\frac{x^{2}-1}{2}\right) = x\left(\frac{3-x^{2}}{2}\right)\\
          11 &= 8x(3-x^{2})\\
          0 &= 8x^{3} -24x + 11\\
          \intertext{By the rational root theorem possible rational zeroes of this cubic are:}
          \pm 1,\pm 11,\pm\frac{1}{2},\pm\frac{11}{2},\pm\frac{1}{4},\pm\frac{11}{4},\pm\frac{1}{8},\pm\frac{11}{8}\\
          \shortintertext{Upon checking (and given in the hint!), we find that \(\frac{1}{2}\) is a root, so we can factorise:}
          0 &= (2x-1)(\_x^{2} + \_x + \_) = 8x^{3} - 24x + 11\\
          \intertext{where the underscores are values to find. The coefficient of \(x^{2}\) must be \(4\) and we can see the constant term must be \(-11\).}
          0 &= (2x-1)(4x^{2} +\_x -11) = 8x^{3} - 24x + 11\\
          \intertext{Now the coefficent of \(x\) must be \(2\) giving}
          0 &= (2x-1)(4x^{2} + 2x -11)\\
          \shortintertext{Solving the quadratic:}
          x &= \frac{1}{2},\frac{-1\pm 3\sqrt{5}}{4} \text{ that is } x = 0.5,\;1.42705\dots,\; -1.92705\dots\\
          \shortintertext{Our possible solutions are:}
          \sin\theta + \cos\theta &= 0.5,\;1.42705\dots,\; -1.92705\dots
        \shortintertext{However, the maximum and minimum values of \(\sin\theta + \cos\theta\) are \(\pm\sqrt{2}\), that is,}
          -\sqrt{2} &\leq \sin\theta + \cos\theta \leq \sqrt{2}\\
       \shortintertext{The two surd roots are out of this range,}
          \frac{-1 + 3\sqrt{5}}{4} > \sqrt{2},&\qquad \frac{-1 - 3\sqrt{5}}{4} < -\sqrt{2}
     \shortintertext{So the only possible solution is}
        \sin\theta + \cos\theta &= \frac{1}{2}
      \end{align*}
      \end{MySolutionBox}
    }{}%conditional output end
    \end{MyInnerBox}


  \end{MyOuterBox}

%%%%%%%%%%%%%%%%%%%%%%%%%%%%%%%%%%%%%%%%%%%%%%%%%%
\end{document}



